\section{Job Description}

A job description (JD) is a document that outlines the purpose of a role, as well as the competencies required.
In the technical landscape, requirements are usually explicitly listed, such as programming languages, frameworks,
tooling, seniority and responsibilities. Hiring managers use this text to decide which skills to prioritise,
and to what extent each of these should influence shortlisting, however, most JDs are written for candidates,
not for machines. As such, there is often prose that is not relevant to the screening process, such as benefits and company culture.

There is currently no industry standard for JDs, hence they differ in terms of clarity and depth.
Some postings detail requirements precisely, while others follow boilerplate text that repeats common terms across roles.
Candidates often form unrealistic expectations due to this lack of clarity and specificity \cite{breaugh1983,earnest2011}. 

Regardless, in most cases, the JD can be seen as a valuable source of information for the screening process, given that
the system is able to separate the requirements from the prose. Many screening systems do not take advantage of this, instead
applying fixed weights and metrics to measure against, even when the JD may vary in terms of stack and seniority level~\cite{gugnani2020,lo2025aihiring}.
A screening system that is not tailored to a specific JD is likely to be generic, and may pick the best ``average'' candidate,
rather than the best candidate for that specific role. Given this, rigid scoring systems may miss strong candidates, and
shortlist candidates that are not a good fit for the role when the screening configuration does not reflect the JD~\cite{fuller2021hiddenworkers}

MERIT must, therefore, treat the JD as a noisy configuration document, separating explicit requirements from redundant
information. It extracts listed technical metrics and suggests relative weights via corpus-aware TF-IDF, which recruiters
can add to, remove or override before scoring.
